\section{Introduction}
\label{sec:intro}
% ScaleQsim order: context, need, representations, constraints, prior work, design, contributions.

Quantum processors with more than 100 qubits have enabled experiments beyond sampling demonstrations, including measurements of observables in interacting physical models~\cite{arute2019quantum,kim2023evidence}.
Observable expectation values characterize quantities such as magnetization and energy in quantum simulation, with applications in physics, chemistry, and materials science~\cite{georgescu2014simulation}.
For example, a 127-qubit processor measured the magnetization of a kicked Ising model at a scale beyond direct state-vector simulation~\cite{kim2023evidence}.
Classical calculations of these observables provide a basis for assessing the accuracy and computational cost of the quantum experiment.
The tractable circuit size depends on the classical representation, its approximation, and the computing resources available for the calculation.

Current processors remain noisy intermediate-scale quantum (NISQ) devices, with errors that affect measured expectation values~\cite{preskill2018quantum}.
Error mitigation estimates an observable from repeated measurements, including measurements at amplified noise levels~\cite{kim2023evidence}.
Classical simulation on high-performance computing (HPC) systems provides an independent reference for evaluating these estimates.
Such references support hardware validation, error-mitigation assessment, and algorithm development, subject to the numerical error of the classical method.
Tensor network and Pauli propagation studies have reproduced observables from the 127-qubit experiment~\cite{tindall2024efficient,begusic2024fast}.
These calculations establish the relevance of classical simulation even when a full state vector exceeds available memory.

Three representations are relevant to the classical simulation of these circuits.
State vector simulation stores all $2^n$ complex amplitudes of an $n$-qubit state, requiring 16~PiB for 50 qubits in double precision~\cite{haner2017petabyte}.
Tensor network simulation represents the state through tensors, with contraction cost determined by the network structure and the required tensor dimensions~\cite{tindall2024efficient}.
Pauli propagation represents an observable as a sum of Pauli strings and propagates it through the circuit in the Heisenberg picture~\cite{begusic2024fast,rudolph2025pauli}.
Its storage and execution costs depend on the number of retained terms as well as the string length and gate sequence.
Current NVIDIA and IBM libraries support this representation~\cite{cupauliprop,qiskitpauliprop}.
Our work focuses on distributed Pauli propagation, where coefficient truncation controls the retained expansion and permits calculations beyond the memory range of full state vectors.

The retained expansion can exceed the capacity of one GPU as the coefficient cutoff decreases.
Terms also accumulate between scheduled truncations, requiring temporary storage beyond the retained expansion.
For the 127-qubit kicked Ising circuit, the maximum population at a window boundary increases from 159 million terms at cutoff $10^{-5}$ to 11.9 billion at $10^{-6}$.
At 40 bytes per term, these populations require about 6.4~GB and 475~GB before execution buffers.
Distributed implementations address this storage requirement on GPU clusters and on CPU systems, including Fugaku runs with $2^{17}$ processes~\cite{cupauliprop,orqa}.
However, the distributed GPU implementations still incur coordination costs at individual gates and require their local expansions to fit device memory.

%   시스템별 나열을 빼고 원인 두 가지(gate마다의 전역 단계, GPU 메모리)와 메모리가 모자라는 이유를 씀. 시스템별 배수는 평가로.
Distributed GPU propagation has two principal constraints.
First, terms with the same Pauli string must be merged before coefficient truncation, including contributions produced on different GPUs.
The distributed mode of \cuPP coordinates this merge after every gate~\cite{cupauliprop}.
The measured fixed cost increases from 5.9~ms per gate on four GPUs to 29.8~ms on 16 GPUs.
Second, device memory must accommodate retained terms, intermediate branches, and merge workspace.
An anticommuting rotation can produce two terms from one input, increasing the expansion between truncations.
Under these requirements, \cuPP retains about 100 million terms per GPU with a 36~GB memory budget.
The GPU implementation of paulistrings-rs also stores its assigned partitions in device memory~\cite{paulistrings}.
Deferred deduplication reduces global merge frequency but increases storage for duplicate strings.
Thus, a distributed design must provide merge locality across multiple gates and accommodate the stored operator beyond the working buffers of individual GPUs.

\begin{table}[t]
\caption{Features of Pauli propagation systems}
\label{tab:prior}
\centering
\small
\begin{tabular}{@{}llccccc@{}}
\toprule
\textbf{System} & \textbf{Method} & \textbf{MN} & \textbf{GPU} & \textbf{LP} & \textbf{RP} & \textbf{HM} \\
\midrule
PauliPropagation.jl~\cite{rudolph2025pauli} & Pauli propagation, one node & & & & & \\
SPD~\cite{spd} & Sparse Pauli dynamics & & \checkmark & & & \\
ORQA~\cite{orqa} & Structured distribution map & \checkmark & & & & \\
MonoProp~\cite{monoprop} & Distributed Pauli propagation & \checkmark & & & & \\
\cuPP~\cite{cupauliprop} & Global merge after every gate & \checkmark & \checkmark & & & \\
paulistrings-rs~\cite{paulistrings} & Fixed linear partition & \checkmark & \checkmark & \checkmark & & \\
\midrule
\textbf{\SectorQ (Ours)} & \textbf{Window-conserved parity sectors} & \checkmark & \checkmark & \checkmark & \checkmark & \checkmark \\
\bottomrule
\end{tabular}
\par\smallskip
\begin{minipage}{\linewidth}\footnotesize
MN denotes multi-node execution, GPU device execution, LP locality-preserving partitions, RP run-time partitions derived from gates and live terms, and HM host storage beyond GPU capacity.
\end{minipage}
\end{table}

Previous systems address arithmetic, distribution, and storage at different levels, as summarized in Table~\ref{tab:prior}.
Single-node simulators use sparse data structures, circuit symmetries, and GPU kernels to reduce propagation cost~\cite{begusic2024fast,rudolph2025pauli,spd}.
ORQA uses a structured distribution map to limit communication partners, and MonoProp supports distributed CPU propagation~\cite{orqa,monoprop}.
The distributed mode of \cuPP combines contributions across GPUs after individual gates~\cite{cupauliprop}.
The linear partition in paulistrings-rs assigns strings through fixed bit parities, with a predictable communication partner when a generator changes a partition label~\cite{paulistrings}.
Commutator-graph analysis also identifies operator subspaces invariant under a set of generators~\cite{west2025graph}.
None of these implementations combines window-conserved rows, row selection from live terms, and host-memory sector execution.

\SectorQ constructs a partition for each window of consecutive gates and selects its rows during execution.
The window generators define conserved parities that assign every contribution to an equal Pauli string to the same sector.
The runtime selects among these parities using a sample of the live terms, then assigns each sector to one node for the window.
This assignment confines merge and truncation to sector execution until the next window boundary.
The host-memory store retains the sectors, and each GPU loads one sector at a time with the workspace required by the propagation library.
An oversized sector is subdivided using further rows valid for its remaining operations.
Operator capacity consequently depends on distributed host storage and the availability of executable sector subdivisions, while term exchange occurs at window boundaries.

% ¶8 "In this paper": ACCEPTED by the user (2026-10-14), keep as it is.
% (ScaleQsim형 대안 "The goal of SectorQ is to 1) 2) 3). To achieve these goals, SectorQ 1) 2) 3)."은 사용자가 고르면 바꿈)
In this paper, we present \SectorQ, a Pauli propagation system for GPU clusters that partitions the terms by the parities that a window of gates conserves.
A rotation gate maps each term to itself and to its product with the generator of the gate, and every parity that the generator leaves unchanged assigns the same value to both terms.
Specifically, \SectorQ introduces 1) window-conserved parity sectors that the system recomputes for every window from the generators of its gates, 2) a run-time choice of the parities and the number of sectors from a sample of the live terms while the GPUs process the current window, and 3) a host-memory sector store from which each GPU loads one sector at a time.
With these sectors, the terms move between nodes only at window boundaries, and the number of terms that \SectorQ holds follows the host memory of all nodes.
Our evaluation on up to 512 NVIDIA A100 GPUs of the Perlmutter supercomputer shows that \SectorQ reduces the simulation time by 4.4-9.7$\times$ on 16 GPUs and increases the largest finished problem by 140$\times$ on 512 GPUs compared with \cuPP, the existing multi-GPU engine.
We open-source \SectorQ at: \url{https://anonymous.4open.science/r/SectorQ}
